\chapter{Related Work}
\label{chap:related-work}

% 5 pages imo

% In this chapter I introduce the literature background for the thesis. I begin with a comprehensive overview of existing and examined SSL training frameworks, focusing on the architectures, training, evaluation, and downstream tasks. Then, I provide a breakdown on adversarial examples, explaining the theory behind them, their universality, current attacks, and pair them with current SSL paradigm.

% \section{SSL Vision Encoders}

% Contrary to classification task, where the goal is to predict a correct class $c\in{1, 2,...,C}$ based on the image. $x\in\mathbb{R}^{C\times H\times W}$, the goal of SSL encoders is more general. Classifiers learn based on labeled data $(x, c)$, under supervised learning paradigm,~\ie the correct prediction is well known, and the sole training objective is to teach the NN to output a specific (correct) class $c$ for each image $x$.

% SSL abandons the luxury of labeled data. Intuitively, most of the data one can obtain via Internet-scale scraping is unlabelled, and the labelling requires costly human labor. Without labels, SSL posits an alternative training objective,~\ie representation learning. Representations $z\in\mathbb{R}^d$ are understood as high-dimensional ($d$ denotes number of dimensions), continuous vectors. After training, the SSL model $f_\theta: \mathbb{R}^{C\times H\times W}\rightarrow\mathbb{R}^d$ (with parameters $\theta$) should output rich and expressive representation $z$ of any image $x$, enabling further downstream use of the representations.

% \subsection{Training Paradigms}

% Notably, what exactly is a "rich and expressive representation $z$ of any image $x$, enabling further downstream use of the representations" is not well-grounded in theory. Moreover, one cannot compute gradient of such an objective. 

% More precise view of the above goal is: representations $z_1$, $z_2$ of two \textit{similar} images $x_1$, $x_2$ should be \textit{closer} than representations $z_1$, $z_3$ of two \textit{dissimilar} images $x_1$, $x_3$. Intuitively, an image of a dog should have a representation that lies close to other images of dogs in the embedding space, and should have a representation that lies far from images depicting circuit boards. 

% \subsection{Similar Images and SSL}

% Defining "similarity" of images is tricky, especially since images can be similar in many ways. The similarity may be structural, compositional, semantic, stylistic, etc. To avoid ambiguities and human-based biases, the default approach is to augment \textit{the same} image multiple times, producing multiple "views". These views are (likely) similar to each other, and the similarity is as broad as imaginable. 

% Representations' "closeness" is easier to precisely formulate. Since embeddings are continuous vectors, the natural approach is to use an abstract \textit{distance} function $\text{dist}: \mathbb{R}^d\times\mathbb{R}^d\rightarrow\mathbb{R}$, for example L2 distance or cosine similarity. 

% Effectively, given a dataset of $N$ unlabeled images $D=\{x_i\},i=1,...,N$ the SSL training objective utilizing multiple views of the same image is as follows:

% \begin{equation}
%     \label{eq:loss_ssl_sim}
%     \theta^*=\underset{\argmin}{\theta}\mathbb{E}_{x\sim D,\mathcal{A_i}\sim{\mathbb{A}},\mathcal{A_j}\sim\mathbb{A}}\left[\text{dist}(\mathcal{A_i}(x),\mathcal{A_j}(x))\right]\text{,}
% \end{equation}
% where $\mathbb{A}={\mathcal{A_1},\mathcal{A_2},...}$ is a set of augmentations $\mathcal{A_i}: \mathbb{R}^{C\times H\times W}\rightarrow\mathbb{R}^{C\times H\times W}$.

% \subsection{The Great Collapse}

% Unfortunately,~\cref{eq:loss_ssl_sim} has a fatal flaw: its minimizer is a network $f_\theta^*: \mathbb{R}^{C\times H\times W}->[0]^d$. To address that HERE SIMSIAM AND DINO ONCE I HAVE INTERNET TO WRITE IT AS GOOD AS IT GETS.

% \subsection{Dissimilar Images and Contrastive Learning}

% With \textit{similarity} defined, I now introduce the first specific SSL training paradigm: Contrastive 

























This chapter expands the research context for the thesis. The goal is not to reproduce every paper in self-supervised learning or adversarial robustness, but to make clear why the evaluation in later chapters is necessary. The central issue is a mismatch between how foundation vision encoders are used and how robust SSL encoders are often evaluated. Foundation encoders are reused through many adaptors, while robust SSL papers frequently report a smaller number of classification-centered metrics.

\section{From Feature Learning to Foundation Encoders}

Representation learning has long separated feature extraction from task-specific prediction. Classical computer vision used hand-designed descriptors, then supervised convolutional networks made learned features dominant. Foundation vision models extend this separation: a large encoder is trained once and then used in many downstream settings. The foundation-model framing emphasizes that a model is not a single application but reusable infrastructure \parencite{bommasani2021opportunities}.

In vision, the encoder is the reusable component. If $f$ maps an image to tokens or a vector, then many downstream heads can be written as $h^t(f(x))$. This decomposition is attractive because pre-training is expensive and labeled data can be scarce. The same pre-trained representation can support a classifier, a segmentation head, a depth head, or retrieval. The stronger the representation, the more tempting it is to evaluate it once and assume that the conclusion generalizes.

This assumption is safe only for some properties. Clean classification performance often correlates with representation quality, but it is not the same as clean dense-prediction quality. Adversarial robustness is even less likely to be a single scalar. A perturbation that changes a class token may differ from one that damages patch tokens. A perturbation that maximizes cross-entropy may differ from one that maximizes depth error. Therefore, once an encoder is described as a foundation model, evaluation must follow the interfaces where the foundation is actually used.

\section{Self-Supervised Learning Objectives}

SSL uses data itself to define supervision. The common intuition is that two semantic-preserving views of the same image should have compatible representations. Contrastive learning implements this by attracting positive pairs and repelling negatives. SimCLR is the canonical modern formulation: it uses strong augmentation, a projection head, and a contrastive loss over many in-batch negatives \parencite{chen2020simple}. The learned representation is then evaluated through a linear classifier.

SimSiam shows that useful representations can be learned without explicit negative pairs \parencite{simsiam}. It uses a Siamese architecture, a predictor on one branch, and a stop-gradient operation on the other. The method is important here because it highlights a broad lesson: SSL representation quality depends on the geometry created by the training objective, not simply on labels. When robustness is later imposed on the representation, it interacts with that geometry.

\DINO uses self-distillation \parencite{caron2021dino}. The student and teacher see different augmented views, and the student learns to match the teacher distribution. The teacher is updated by an exponential moving average. This produces ViT features with strong emergent structure. In the context of this thesis, \DINO is important not only because it performs well, but because its token structure naturally supports dense tasks. The same model can supply a global representation for classification and patch tokens for segmentation or depth.

Masked autoencoders provide another SSL family \parencite{mae}. Instead of matching views or contrasting pairs, they reconstruct missing image patches. They demonstrate that pretext objectives can train very strong ViT encoders. This thesis does not evaluate MAE directly, but its presence in the resource bibliography matters: modern SSL is diverse, and a robust evaluation protocol should not be tied to one pretext task.

\DINOv further strengthens the foundation-model motivation \parencite{oquab2024dinov2}. It is explicitly presented as a general-purpose visual representation useful for many downstream tasks. It also uses dense outputs and large-scale training to improve transfer. If a model family is deployed because it supports many tasks, then robustness should be measured across those tasks. A classification-only robustness number is not enough evidence for a \DINOv-style use case.

\section{Downstream Evaluation Protocols}

Linear probing is the standard SSL protocol. The encoder is frozen, a linear classifier is trained, and test accuracy is reported. This protocol is useful because it isolates the linear separability of the representation. It is also cheap compared with full fine-tuning. Many robust SSL papers therefore report clean and adversarial linear-probe accuracy.

However, linear probing is a narrow interface. Classification compresses the image into one label. Semantic segmentation predicts a label for every pixel. The segmentation head must preserve spatial information, combine local and global context, and upsample token features. A representation can be good for global category prediction while being weaker for boundary localization or small objects. Fully convolutional networks and UPerNet represent the dense prediction tradition in the bibliography \parencite{long2017segmentation,xiao2018unified}.

Depth estimation is different again. The output is continuous, not categorical. Errors are geometric and spatially structured. Losses often mix pixel-level depth differences with gradient consistency \parencite{godard2019depthestimation,MegaDepthLi18,AdaBins}. A perturbation that causes a small classification error may cause a large depth distortion, or vice versa. This makes depth estimation a useful stress test for the idea that representation robustness transfers.

The thesis uses these three tasks because they form a compact but meaningful set. Classification tests global semantics. Segmentation tests spatial semantics. Depth tests geometry. Together they cover several ways in which a foundation encoder can be reused. They are still not exhaustive: detection, optical flow, video, medical imaging, and multimodal grounding remain outside scope. But the set is broad enough to show that a single classification metric is inadequate.

\section{Adversarial Examples and PGD}

Adversarial examples show that small perturbations can cause large prediction changes \parencite{biggio2013evasion,szegedy2014intriguing}. In modern deep learning, the standard white-box setting assumes the attacker knows the model and differentiates through the loss. Projected gradient descent is a widely used first-order method for this purpose \parencite{madry2018towards}. Starting from a point inside the perturbation ball, it repeatedly moves in the sign of the loss gradient and projects back to the feasible set.

For supervised classification, the objective is usually cross-entropy. For dense prediction, the objective must reflect the dense output. Segmentation attacks sum or weight pixel-wise cross-entropy. Depth attacks maximize a depth loss. This thesis treats those objectives as different interfaces. They are not merely implementation details; they define what the attacker wants to break.

The distinction matters because adversarial robustness is often local to the training objective. A model adversarially trained against one norm may still fail under another norm. Supervised work on multiple perturbations makes this point directly \parencite{tramer2019adversarial}. The same reasoning applies to multiple tasks. A model robust to representation-distance perturbations may not be robust to segmentation-loss perturbations.

\section{Adversarial Robustness in SSL}

Robust SSL adapts adversarial ideas to the unlabeled setting. Since labels are absent during pre-training, methods define adversarial examples through representations, contrastive losses, or consistency objectives. RoCL uses adversarial views in contrastive learning \parencite{kim2020adversarial}. ACL encourages consistency between clean and adversarial examples \parencite{jiang2020robust}. AdvCL studies when contrastive learning preserves adversarial robustness through fine-tuning \parencite{fan2021does}. CLAE builds contrastive examples using adversarial perturbations \parencite{ho2020CLAE}. DYNACL reconsiders data augmentation and robustness trade-offs \parencite{luo2023rethinking}.

These works demonstrate progress, but they also show why evaluation choices matter. If adversarial examples are generated by maximizing a representation or contrastive objective, then the resulting defense is naturally aligned with that objective. It may still help downstream tasks, but the help is empirical rather than guaranteed. The downstream head can expose gradients and error directions that the SSL objective never constrained.

\DeACL is especially relevant because it decouples pre-training from robust fine-tuning \parencite{zhang2022decoupled}. This is practical for foundation models. Large SSL pre-training can remain unchanged, and a robust encoder can be obtained afterward by distillation. The robust encoder tries to preserve clean representations while aligning adversarial and clean inputs. The evaluated question is whether this representation alignment is enough for dense downstream robustness.

\section{Feature Robustness and Transfer}

A key conceptual debate concerns what adversarial features are. Some work argues that adversarial examples exploit predictive but non-robust features \parencite{ilyas2019adversarialbugs}. Later work questions whether adversarial features from supervised models transfer as useful SSL features \parencite{li2023adversarialnotbugs}. This debate matters because SSL is often claimed to learn different, sometimes more semantically meaningful, representations than supervised training.

The thesis does not attempt to settle the feature debate theoretically. Instead, it asks an operational question. Suppose a self-supervised encoder is clean-useful across tasks. Suppose a defense makes the representation less sensitive to representation-level adversarial perturbations. Do downstream users obtain robustness? The answer in the supplied experiments is mostly no. This operational conclusion is important even if the theoretical explanation remains open.

\section{Evaluation Gaps Addressed by This Thesis}

The related work leaves three gaps. First, robust SSL evaluation is too often centered on classification. Second, representation attacks are sometimes treated as general proxies for downstream attacks. Third, dense prediction tasks are underused in robustness claims about reusable encoders.

This thesis addresses these gaps by reporting a task-and-attack matrix. Rows are datasets and tasks; columns are clean, \EmbedAttack, and task-specific attack conditions. The matrix does not solve robust SSL, but it prevents a misleading scalar summary. It shows exactly where \DeACL helps and where it does not.

The result is also a recommendation for future papers. If an encoder is described as a foundation vision model, robust evaluation should include both sparse and dense tasks. If a defense optimizes a representation objective, the paper should report representation attacks and downstream attacks separately. If clean performance drops, that trade-off should be stated for every downstream task. These reporting norms follow directly from the related-work gap.
